% =====================================================================
%  CHAPTER 9: 儿童少年慢性病
% =====================================================================
\chapter{儿童少年慢性病}
\setcounter{chapter}{9}
\chapterintro{
掌握儿童少年常见慢性病的类型；掌握儿童糖尿病（主要为1型）的流行病学特征和防治；掌握儿童高血压的定义、危险因素及预防；熟悉动脉粥样硬化的早期危险因素及预防策略；理解生命早期干预对成人慢性病预防的重要意义。
}

% ----- 9.1 概述 -----
\section{儿童少年慢性病概述}

\begin{defbox}
\textbf{儿童少年慢性病}：在儿童少年期起病或具有早期危险因素的慢性非传染性疾病（NCDs）。随着疾病谱从传染病向慢性病转变，儿童慢性病已成为全球重要的公共卫生问题。

\textbf{常见类型：}
\begin{itemize}[leftmargin=*]
    \item \imp{糖尿病}（主要为1型糖尿病，2型糖尿病在肥胖儿童中增多）
    \item \imp{高血压}
    \item \imp{动脉粥样硬化}（病理过程可从儿童期开始）
    \item \imp{哮喘}
    \item \imp{肥胖}（既是常见病，也是其他慢性病的重要危险因素）
\end{itemize}

\textbf{生命早期起源假说（Barker假说/DOAD假说）}：胎儿期和生命早期的营养不良及不良暴露，可通过"编程"机制增加成年期高血压、糖尿病、冠心病等慢性病的发生风险。这强调了儿童期慢性病防控对终身健康的重要意义。
\end{defbox}

% ----- 9.2 糖尿病 -----
\section{儿童糖尿病}

\subsection{概述}

\begin{keybox}
\textbf{儿童糖尿病的类型：}

\begin{enumerate}[leftmargin=*]
    \item \imp{1型糖尿病（T1DM）}——儿童期最常见的糖尿病类型（占90\%以上），为自身免疫性胰岛$\beta$细胞破坏导致胰岛素绝对缺乏。发病高峰为10--14岁（青春期）
    \item \imp{2型糖尿病（T2DM）}——近年来随儿童肥胖率上升而增多，以胰岛素抵抗为特征。多见于肥胖、有家族史的青少年
    \item \imp{特殊类型糖尿病}——如青少年发病的成人型糖尿病（MODY），由单基因突变引起
\end{enumerate}

\textbf{1型糖尿病的临床特点：}
\begin{itemize}[leftmargin=*]
    \item 起病急，"三多一少"（多饮、多尿、多食、体重减少）症状典型
    \item 需要终身胰岛素治疗
    \item 急性并发症：糖尿病酮症酸中毒（DKA）是最常见的急性并发症，可危及生命
    \item 慢性并发症：微血管病变（视网膜病变、肾病、神经病变）
\end{itemize}
\end{keybox}

\subsection{防治策略}

\begin{enumerate}[leftmargin=*]
    \item \imp{早期识别与及时诊断}——警惕"三多一少"典型症状
    \item \imp{规范胰岛素治疗}——模拟生理胰岛素分泌模式
    \item \imp{血糖监测}——自我血糖监测和糖化血红蛋白（HbA1c）定期检测
    \item \imp{医学营养治疗}——个体化饮食方案
    \item \imp{规律运动}——有助于血糖控制
    \item \imp{心理社会支持}——长期慢性病管理的心理负担重大
    \item \imp{学校管理}——学校应提供支持性环境，协助血糖监测和胰岛素使用
\end{enumerate}

% ----- 9.3 高血压 -----
\section{儿童高血压}

\subsection{定义与诊断}

\begin{defbox}
\textbf{儿童高血压}：儿童血压超过同性别、年龄和身高别的第95百分位值。

\textbf{诊断标准（中国）：}
\begin{itemize}[leftmargin=*]
    \item 正常血压：收缩压和舒张压 $<$ P90
    \item 正常高值血压（高血压前期）：P90 $\leq$ 血压 $<$ P95（或$\geq$120/80mmHg）
    \item 高血压：血压 $\geq$ P95。分两期：
    \begin{itemize}
        \item 1期高血压：P95 $\leq$ 血压 $<$ P99 + 5mmHg
        \item 2期高血压：血压 $\geq$ P99 + 5mmHg
    \end{itemize}
\end{itemize}

\textbf{原发性高血压 vs 继发性高血压}：儿童期高血压以继发性为主（肾脏疾病是最常见继发原因）；随肥胖率上升，原发性高血压比例逐渐增加。
\end{defbox}

\subsection{危险因素}

\begin{keybox}
\textbf{儿童高血压的主要危险因素：}
\begin{enumerate}[leftmargin=*]
    \item \imp{超重和肥胖}——\keyword{最重要的可改变危险因素}
    \item \imp{家族史}——有高血压家族史的儿童发病风险增高
    \item \imp{膳食因素}——高盐摄入、低钾低钙饮食
    \item \imp{体力活动不足}
    \item \imp{静态行为过多}
    \item \imp{低出生体重}——生命早期不良暴露与后续高血压风险相关
    \item \imp{睡眠问题}——睡眠不足或睡眠呼吸暂停
\end{enumerate}
\end{keybox}

\subsection{预防与控制}

\begin{enumerate}[leftmargin=*]
    \item \imp{定期血压监测}——3岁以上儿童每年至少测量一次血压
    \item \imp{生活方式干预}（一线治疗）：
    \begin{itemize}
        \item 控制体重（最重要的干预措施）
        \item 减少钠盐摄入
        \item 增加蔬菜水果摄入
        \item 增加体力活动（每天$\geq$60min中高强度运动）
        \item 减少屏幕时间
    \end{itemize}
    \item \imp{药物治疗}——生活方式干预无效或2期高血压时考虑
\end{enumerate}

% ----- 9.4 动脉粥样硬化 -----
\section{动脉粥样硬化早期预防}

\begin{defbox}
\textbf{动脉粥样硬化（Atherosclerosis）}：是一种始于儿童期的慢性进行性血管疾病过程。尸检研究（如Bogalusa心脏研究、PDAY研究）已证实，动脉粥样硬化的早期病变（脂纹）可在儿童期甚至婴幼儿期出现，且病变程度与心血管危险因素的水平密切相关。

\textbf{儿童期可改变的心血管危险因素：}
\begin{enumerate}[leftmargin=*]
    \item \imp{血脂异常}——高LDL-C、低HDL-C
    \item \imp{肥胖}（特别是中心性肥胖）
    \item \imp{高血压}
    \item \imp{吸烟和二手烟暴露}
    \item \imp{缺乏体力活动}
    \item \imp{不健康膳食}（高饱和脂肪、高反式脂肪、高糖饮食）
\end{enumerate}

\textbf{预防策略的关键}：\keyword{生命早期干预是预防成人动脉粥样硬化性心血管疾病的最有效策略。}
\end{defbox}

% ----- 复习题 -----
\section{复习题}

\begin{enumerate}[leftmargin=*, itemsep=8pt]
    \item \textbf{【名词解释】}1型糖尿病；儿童高血压；动脉粥样硬化
    \item \textbf{【名词解释】}生命早期起源假说（Barker假说）
    \item \textbf{【简答题】}简述儿童1型糖尿病的临床特点和防治原则。
    \item \textbf{【简答题】}简述儿童高血压的诊断标准及主要危险因素。
    \item \textbf{【简答题】}为什么动脉粥样硬化的预防应从儿童期开始？
\end{enumerate}

\subsection*{参考答案要点}

\begin{enumerate}[leftmargin=*, itemsep=5pt]
    \item \textbf{1型糖尿病}：自身免疫性胰岛$\beta$细胞破坏导致胰岛素绝对缺乏，是儿童最常见糖尿病类型（90\%以上）。\textbf{儿童高血压}：血压超过同性别、年龄、身高别的P95。\textbf{动脉粥样硬化}：慢性进行性血管病变，早期病变（脂纹）可在儿童期出现。

    \item \textbf{生命早期起源假说}：胎儿期和生命早期不良暴露可通过"编程"机制增加成年期慢性病风险，强调生命早期干预的极端重要性。

    \item 特点：起病急、"三多一少"、需终身胰岛素治疗、DKA是最常见急性并发症。防治：早期识别诊断、规范胰岛素治疗、血糖监测、医学营养治疗、规律运动、心理支持、学校管理。

    \item 诊断：正常（$<$P90）→正常高值（P90--P95）→高血压（$\geq$P95，分1期和2期）。危险因素：超重肥胖（最重要）、家族史、高盐饮食、体力活动不足、静态行为多、低出生体重。

    \item Bogalusa心脏研究和PDAY研究证实动脉粥样硬化早期病变可在儿童期出现；病变程度与心血管危险因素水平相关；儿童期危险因素可追踪至成人；早期干预比成人期干预更有效。生命早期是预防的最关键窗口期。
\end{enumerate}

\newpage
